\documentclass[11pt]{article}
\usepackage[margin=1.1in]{geometry}
\usepackage{amsmath,amssymb,amsthm}
\usepackage{booktabs}
\usepackage{enumitem}
\usepackage[round]{natbib}
\setcitestyle{aysep={}}
\usepackage{xcolor}
\definecolor{linknavy}{RGB}{25,55,120}
\usepackage[colorlinks=true,linkcolor=linknavy,citecolor=linknavy,urlcolor=linknavy,
pdftitle={Core stability does not imply egalitarian stability},
pdfauthor={Maxwell Flanagan}]{hyperref}

\newtheorem{theorem}{Theorem}[section]
\newtheorem{lemma}[theorem]{Lemma}
\newtheorem{proposition}[theorem]{Proposition}
\newtheorem{corollary}[theorem]{Corollary}
\newtheorem{definition}[theorem]{Definition}
\newtheorem{example}[theorem]{Example}
\newtheorem{remark}[theorem]{Remark}

\newcommand{\R}{\mathbb{R}}
\newcommand{\core}{\operatorname{Core}}

\title{Core stability does not imply egalitarian stability:\\
a sharp six-player counterexample}

\author{Maxwell Flanagan\\[2pt]
\normalsize Corresponding author: \texttt{maxwell.flanagan@gmail.com}%
\thanks{A portable package contains the witness data and exact-arithmetic verification code.}}

\date{Specialist-review draft, September 5, 2026}

\begin{document}
\maketitle

\begin{abstract}
\citet{Dietzenbacher2021} asked whether every transferable-utility game whose
core is a von Neumann--Morgenstern stable set is egalitarian stable, meaning
that the procedural egalitarian solution selects a core allocation. We answer
no by applying a player-specific additive perturbation to a classical
six-player stable-core game. The perturbation translates the core and
preserves domination, exactness, monotonicity, and superadditivity, but its
egalitarian claims sum to $13/2$, exceeding the grand-coalition worth $6$.
We give a self-contained stability proof: five tight pairs reduce every
imputation not dominated by the core to a one-parameter discrepancy, and one
of two tight triples eliminates either sign. The player bound is sharp,
because core stability and core largeness coincide for at most five players,
and large-core games are egalitarian stable. The construction also shows that
egalitarian stability is not invariant under player-specific additive
transformations. In fact, a one-parameter additive ray is egalitarian stable
only at its unperturbed endpoint, with three explicit affine claim regimes.
The exact-game half of Dietzenbacher's question had already
been exhibited by \citet{Mauri2017}.
\end{abstract}

\medskip
\noindent\textbf{Keywords:} transferable-utility games; stable core; stable
set; procedural egalitarian solution; exact games.

\section{Introduction}\label{sec:intro}

The procedural egalitarian solution extends the constrained egalitarian
solution of \citet{DuttaRay89} beyond convex games \citep{DBH2017}. Its
natural domain is the class of
\emph{egalitarian stable} games: those on which the procedure selects from the
core \citep{Dietzenbacher2021}. This class contains all games with a large core.
Two familiar classes therefore supplied unresolved boundary questions.
\citet{Dietzenbacher2021} asked whether every exact game is egalitarian stable
and whether the same is true whenever the core is a von
Neumann--Morgenstern stable set.

The two halves have different prior-art histories.
\citet[][Example~5, pp.~40--41]{Mauri2017} applied the relevant maximum-average procedure
to the five-player exact game of \citet{Biswas1999}, obtained the zero vector,
and observed that it lies outside the core. In the later terminology of
\citet{Dietzenbacher2021}, that example is not egalitarian stable. We recall
this point in Section~\ref{sec:exact-prior} but claim no novelty for the
exact-game half.

For the stable-core half, we start from a classical six-player orthogonal
game. Six-player stable-but-not-large examples go back at least to
\citet{Biswas2001}; the precise segment-core game used here is recorded by
\citet[][Example~4.3]{Shellshear2009}. We add the nonuniform additive vector
$c=(1,0,0,0,1,0)$. After a permutation of the classical game, the perturbed
game is generated by
\[
 A=(1,1,0,1,3,0),\qquad B=(2,0,1,0,1,2),\qquad
 v(S)=\min\{A(S),B(S)\}.
\]
Its core is exactly the segment $[A,B]$, and its egalitarian claims are
\[
 \widehat\gamma
   =\left(\tfrac32,\tfrac12,1,\tfrac12,\tfrac32,\tfrac32\right),
 \qquad
 \widehat\gamma(N)=\tfrac{13}{2}>6=v(N).
\]
Thus the perturbed game is not egalitarian stable. Additive transformations
translate imputations and cores and preserve the dominance relation, so its
stable core is inherited from the classical game. We also give a short,
self-contained proof in the perturbed coordinates: five tight pairs and two
tight triples suffice to dominate every imputation outside the segment. The
new ingredient is therefore not the underlying core geometry, but the fact
that this particular perturbation preserves that geometry while destroying
egalitarian stability.

The result is sharp. \citet{Estevez2012} proved that core stability and core
largeness coincide for games with at most five players, while large-core games
are egalitarian stable \citep{Dietzenbacher2021}. Hence six is the first
possible number of players. The example lies in the class of exact, monotone, superadditive games.

Our contribution is the explicit stable-core/PES separation, its sharp player
count, and its behavior along one additive ray. Recent work studies other
boundaries of egalitarian selection: population monotonicity and egalitarian
allocation schemes \citep{DietzenbacherDogan2025}, and consistency under
core-bound reduced games \citep{GongDietzenbacher2026}. These results do not by
themselves establish the stable-core inclusion considered here. A dated,
claim-level source comparison accompanies this review package. A complete
specialist priority assessment, including comparison with the final 2025
article beyond its accessible working-paper version, remains pending.

\section{Preliminaries}\label{sec:prelim}

Let $N=\{1,\dots,n\}$. A transferable-utility game is a function
$v\colon 2^N\to\R$ with $v(\emptyset)=0$. For $x\in\R^N$ and $S\subseteq N$,
write $x(S)=\sum_{i\in S}x_i$. The core is
\[
 \core(v)=
 \{x\in\R^N:x(N)=v(N),\ x(S)\ge v(S)\text{ for every }S\subseteq N\}.
\]
An \emph{imputation} is an efficient and individually rational vector:
$x(N)=v(N)$ and $x_i\ge v(\{i\})$ for every $i$. Given imputations $x,y$,
$x$ \emph{dominates} $y$ via a nonempty coalition $S$ if
\[
 x_i>y_i\quad(i\in S),\qquad x(S)\le v(S).
\]
The core is internally stable. It is a von Neumann--Morgenstern
\emph{stable set} precisely when it is externally stable: every imputation
outside the core is dominated by a core element \citep{vNM1944}. We then say
that $v$ has a \emph{stable core}, and require this core to be nonempty.
This nonemptiness condition excludes vacuous external stability when the
imputation set is empty.

Following \citet{Schmeidler1972}, a game is \emph{exact} if every coalition
constraint is attained in the core:
\[
 v(S)=\min_{x\in\core(v)}x(S)\qquad(S\subseteq N).
\]
The core is \emph{large} if every vector in the upper core,
$\{y:y(S)\ge v(S)\ \forall S\}$, lies above some core point
\citep{Sharkey1982}. Every large-core game has a stable core and is
egalitarian stable \citep{Sharkey1982,Dietzenbacher2021}.

\subsection{The procedural egalitarian claims}

We recall only the part of the procedural egalitarian solution needed below
\citep{DBH2017,Dietzenbacher2021}. Let $P^0=\emptyset$. At round $k$, claims
$\gamma_i$ have already been fixed for $i\in P^{k-1}$. For each nonempty
$S\subseteq N$, define
\[
 \chi_i^k(S)=
 \begin{cases}
  \gamma_i,
    &i\in S\cap P^{k-1},\\[3pt]
  \displaystyle
  \frac{v(S)-\sum_{j\in S\cap P^{k-1}}\gamma_j}
       {|S\setminus P^{k-1}|},
    &i\in S\setminus P^{k-1},
 \end{cases}
\]
with the first line applying to every member when $S\subseteq P^{k-1}$.
Let $m_i^k=\max_{S\ni i}\chi_i^k(S)$. The round-$k$ admissible collection is
\[
 {\cal A}^k=
 \left\{S\ne\emptyset:
   \sum_{i\in S}\chi_i^k(S)=v(S)
   \text{ and }\chi_i^k(S)=m_i^k\text{ for all }i\in S
 \right\}.
\]
The players in $\bigcup_{S\in{\cal A}^k}S$ are added to $P^k$, and each newly
fixed player $i$ receives claim $\gamma_i=m_i^k$. The procedure terminates
with the \emph{egalitarian claims} vector $\widehat\gamma$. A game is
\emph{egalitarian stable} (or PES-stable) if the grand coalition belongs to
the terminal admissible collection. On exactly this class, the procedural
egalitarian solution selects a core allocation \citep{Dietzenbacher2021}.

\begin{lemma}\label{lem:oversum}
If $\widehat\gamma(N)>v(N)$, then $v$ is not egalitarian stable.
\end{lemma}

\begin{proof}
Suppose instead that $N$ belongs to the terminal admissible collection. Every
player fixed before that round has $\chi_i(N)=\widehat\gamma_i$, while every
player newly fixed in that round receives the value $\chi_i(N)$ because that
player belongs to $N$. Hence $\chi_i(N)=\widehat\gamma_i$ for every $i$ under
this supposition. Admissibility of $N$ also requires
$\sum_i\chi_i(N)=v(N)$, and therefore $\widehat\gamma(N)=v(N)$, contrary to
$\widehat\gamma(N)>v(N)$.
\end{proof}

\subsection{A lower-envelope generator}

Both the prior-work observation and the new witness use the following
standard finite form of the lower-envelope description of exact games. It is
also the elementary finite version of the min-of-additive representation used
in the totally balanced and stable-core literature; see, for example,
\citet{Shellshear2009}.

\begin{lemma}[finite lower envelopes]\label{lem:generator}
Let $Q\subset\R^N$ be a nonempty finite set whose vectors have common sum
$\tau$, and define
\[
 v(S)=\min_{q\in Q}q(S).
\]
Then $v(N)=\tau$, every $q\in Q$ lies in $\core(v)$, and $v$ is exact. If all
coordinates of the generators are nonnegative, then $v$ is monotone and
superadditive.
\end{lemma}

\begin{proof}
For each $q\in Q$, $q(N)=\tau=v(N)$ and
$q(S)\ge\min_{p\in Q}p(S)=v(S)$, so $q\in\core(v)$. Given $S$, choose a
generator $q^S$ attaining the minimum. Then
$q^S(S)=v(S)$, proving exactness.

If the generators are nonnegative, each $q(S)$ is monotone in $S$, hence so
is their pointwise minimum. If $S$ and $T$ are disjoint, then every $q\in Q$
satisfies
\[
 q(S\cup T)=q(S)+q(T)\ge v(S)+v(T).
\]
Taking the minimum over $q$ proves superadditivity.
\end{proof}


\section{The exact-game half is prior work}
\label{sec:exact-prior}

The exact-game half of \citet{Dietzenbacher2021} requires no new
counterexample. For five players, \citet[][Example~5, pp.~40--41]{Mauri2017} considers
the exact game of \citet{Biswas1999}
\[
 a=(1,-1,-1,0,0),\qquad b=(0,1,0,-1,-1),\qquad
 u(S)=\min\{a(S),b(S)\}.
\]
Both generators sum to $-1$, while the maximum coalition average is $0$.
Mauri Masdeu computes the maximizing collection
\[
 \{1\},\quad\{1,2\},\quad\{1,3\},\quad\{1,2,4\},\quad\{1,2,5\},
\]
which covers all players, and obtains the vector $(0,0,0,0,0)$ outside the
core. In the terminology recalled in Section~\ref{sec:prelim}, these are
exactly the first-round admissible coalitions and terminal egalitarian claims.
Since $\widehat\gamma(N)=0>-1=u(N)$, Lemma~\ref{lem:oversum} says that the
game is not egalitarian stable. Thus both the construction and the decisive
egalitarian computation predate the 2021 question; this paper makes no
novelty claim for that half.


\section{A perturbation of a classical stable-core game}\label{sec:witness}

We first record what an additive transformation preserves. For
$c\in\R^N$, write $c(S)=\sum_{i\in S}c_i$.

\begin{lemma}[additive transformations]\label{lem:additive}
If $v(S)=w(S)+c(S)$ for every $S\subseteq N$, then
\[
 I(v)=I(w)+c,\qquad \core(v)=\core(w)+c,
\]
where $I$ denotes the imputation set. Moreover, $v$ is exact if and only if
$w$ is exact, and $v$ has a stable core if and only if $w$ does.
\end{lemma}

\begin{proof}
For every $x\in\R^N$ and $S\subseteq N$,
\[
 (x+c)(S)\ge v(S)\quad\Longleftrightarrow\quad x(S)\ge w(S),
\]
and efficiency and individual rationality transform in the same way. This
proves the two translations and the equivalence of exactness. Finally,
$x+c$ dominates $y+c$ through $S$ in $v$ if and only if $x$ dominates $y$
through $S$ in $w$: the strict coordinate comparisons cancel $c$, as does
the affordability inequality. External stability is therefore preserved.
\end{proof}

Let $N=\{1,\dots,6\}$ and consider
\[
 U=(0,1,0,1,2,0),\qquad V=(1,0,1,0,0,2),\qquad
 w(S)=\min\{U(S),V(S)\}.
\]
Under the coordinate order $(5,2,4,6,1,3)$, its generators become
\[
 (2,1,1,0,0,0)\quad\text{and}\quad(0,0,0,2,1,1).
\]
This is the classical exact stable-core orthogonal game recorded by
\citet[][Example~4.3]{Shellshear2009}; its core is the segment between the
two generators. We use the nonuniform additive vector
$c=(1,0,0,0,1,0)$. The perturbed game $v=w+c$ has generators
\[
 A=U+c=(1,1,0,1,3,0),\qquad
 B=V+c=(2,0,1,0,1,2),
\]
and hence
\begin{equation}\label{eq:game}
 v(S)=\min\{A(S),B(S)\}.
\end{equation}
Both generators sum to $6$ and are nonnegative. By
Lemma~\ref{lem:generator}, $v$ is exact, monotone, and superadditive.
Lemma~\ref{lem:additive} and the cited classical example already imply that
its core is stable. The next two sections give a self-contained segment
description and stability proof in the perturbed coordinates.

\subsection{The core is a segment}

Put
\begin{equation}\label{eq:xt}
 x_t=(1-t)A+tB
     =(1+t,1-t,t,1-t,3-2t,2t),
 \qquad 0\le t\le1.
\end{equation}

\begin{proposition}\label{prop:core}
The core of $v$ is exactly $\{x_t:0\le t\le1\}$.
\end{proposition}

\begin{proof}
Lemma~\ref{lem:generator} puts $A$ and $B$ in the core, so convexity of the
core gives $x_t\in\core(v)$ for $0\le t\le1$.

Conversely, let $y\in\core(v)$. The three disjoint pair constraints
\[
 y_1+y_2\ge2,\qquad y_3+y_4\ge1,\qquad y_5+y_6\ge3
\]
sum to $y(N)\ge6=v(N)$, so efficiency makes all three equalities. Individual
rationality then permits the parametrization
\[
 y=(1+t,1-t,s,1-s,3-u,u),
 \qquad t,s\in[0,1],\quad u\in[0,2].
\]
The fully tight pair constraints for $\{2,3\}$ and $\{1,4\}$ give,
respectively,
\[
 1-t+s\ge1,\qquad 2+t-s\ge2,
\]
and hence $s=t$. Finally, the fully tight triples $\{1,3,5\}$ and
$\{2,4,6\}$ have worths $4$ and $2$. Their core inequalities become
\[
 4+2t-u\ge4,\qquad 2-2t+u\ge2,
\]
so $u=2t$. Therefore $y=x_t$.
\end{proof}

\subsection{The egalitarian procedure leaves the core}

The full coalition function is already specified by
Equation~\eqref{eq:game}; its $63$ nonempty values can therefore be checked
without a separate table. Direct evaluation of the procedure gives the
two-round trace in Table~\ref{tab:trace}. In the first round, the maximum
equal shares for players $1,\dots,6$ are
\[
 \left(\tfrac32,\tfrac54,\tfrac43,\tfrac54,\tfrac32,\tfrac32\right).
\]
Only $\{1,5\}$ and $\{5,6\}$ simultaneously maximize the shares of all their
members, so players $1,5,6$ are fixed at $3/2$. The second round fixes the
remaining players.

\begin{table}[ht]
\centering
\small
\caption{Procedural egalitarian trace for the perturbed six-player game; a
digit string denotes the corresponding coalition.}
\label{tab:trace}
\begin{tabular}{@{}c p{3.55in} l l@{}}
\toprule
round & admissible collection ${\cal A}^k$
      & newly fixed & new claims\\
\midrule
1 & $15,\ 56$
  & $1,5,6$ & $\tfrac32,\tfrac32,\tfrac32$\\[2pt]
2 & $12,\ 14,\ 15,\ 135,\ 56,\ 1256,\ 1456$
  & $2,3,4$ & $\tfrac12,1,\tfrac12$\\
\bottomrule
\end{tabular}
\end{table}

Thus the terminal claims are
\begin{equation}\label{eq:claims}
 \widehat\gamma
 =\left(\tfrac32,\tfrac12,1,\tfrac12,\tfrac32,\tfrac32\right),
 \qquad
 \widehat\gamma(N)=\tfrac{13}{2}>6=v(N).
\end{equation}
By Lemma~\ref{lem:oversum}, $v$ is not egalitarian stable. The verification package
recomputes the trace from Equation~\eqref{eq:game} using exact rational
arithmetic.

\section{The core is externally stable}\label{sec:stability}

The segment in Equation~\eqref{eq:xt} has five useful two-player
projections:
\begin{center}
\begin{tabular}{@{}c c c@{}}
\toprule
pair $P$ & $v(P)$ & $x_t|_P$\\
\midrule
$\{1,2\}$ & $2$ & $(1+t,1-t)$\\
$\{3,4\}$ & $1$ & $(t,1-t)$\\
$\{5,6\}$ & $3$ & $(3-2t,2t)$\\
$\{2,3\}$ & $1$ & $(1-t,t)$\\
$\{1,4\}$ & $2$ & $(1+t,1-t)$\\
\bottomrule
\end{tabular}
\end{center}
In each row, as $t$ ranges over $[0,1]$, the projection sweeps the entire
two-player imputation segment for the corresponding subgame.

\begin{lemma}[pair domination]\label{lem:pair}
Let $P$ be one of the five pairs in the table. If an imputation $y$ satisfies
$y(P)<v(P)$, then some $x_t$ dominates $y$ via $P$.
\end{lemma}

\begin{proof}
Write $P=\{i,j\}$ and
$\varepsilon=v(P)-y(P)>0$. The vector
\[
 z=(y_i+\varepsilon/2,\ y_j+\varepsilon/2)
\]
is individually rational for the two-player subgame and sums to $v(P)$.
Hence it lies on the full segment swept by $x_t|_P$, so
$z=x_t|_P$ for some $t\in[0,1]$. Both members of $P$ strictly gain, and
$x_t(P)=v(P)$; thus $x_t$ dominates $y$ via $P$.
\end{proof}

\begin{theorem}\label{thm:main}
The game in Equation~\eqref{eq:game} is an exact, monotone, superadditive
six-player game with a stable core that is not egalitarian stable.
\end{theorem}

\begin{proof}
Exactness, monotonicity, and superadditivity follow from
Lemma~\ref{lem:generator}; non-egalitarian-stability follows from
Equation~\eqref{eq:claims}. It remains to prove external stability.

Let $y$ be an imputation that is not dominated by any core element.
Lemma~\ref{lem:pair} implies
\[
 y_1+y_2\ge2,\quad
 y_3+y_4\ge1,\quad
 y_5+y_6\ge3,
\qquad
 y_2+y_3\ge1,\quad
 y_1+y_4\ge2.
\]
Efficiency makes the first three inequalities equalities. As in the proof of
Proposition~\ref{prop:core}, individual rationality and the last two
inequalities then give
\[
 y=(1+t,1-t,t,1-t,3-u,u),
 \qquad t\in[0,1],\quad u\in[0,2].
\]

Suppose first that $u>2t$. Choose $r$ strictly between $t$ and $u/2$.
The bounds on $t$ and $u$ ensure $r\in[0,1]$, and
\[
 x_r|_{\{1,3,5\}}=(1+r,r,3-2r)
   >(1+t,t,3-u)=y|_{\{1,3,5\}}.
\]
Since $x_r(\{1,3,5\})=4=v(\{1,3,5\})$, $x_r$ dominates $y$ through that
triple, a contradiction.

If $u<2t$, choose $r$ strictly between $u/2$ and $t$. Then
\[
 x_r|_{\{2,4,6\}}=(1-r,1-r,2r)
   >(1-t,1-t,u)=y|_{\{2,4,6\}},
\]
and $x_r(\{2,4,6\})=2=v(\{2,4,6\})$, again a contradiction. Consequently
$u=2t$, so $y=x_t\in\core(v)$ by Proposition~\ref{prop:core}. Every
imputation outside the core is therefore dominated by a core element.
\end{proof}

\begin{corollary}[sharpness]\label{cor:sharp}
Six is the smallest number of players for which a game can have a stable core
without being egalitarian stable.
\end{corollary}

\begin{proof}
Theorem~\ref{thm:main} supplies a six-player game. For at most five players,
core stability is equivalent to core largeness \citep{Estevez2012}; the
unrestricted TU-game statement is Theorem~4.1 in the author's 2011 working
paper. Every large-core game is egalitarian stable by
\citet[Theorem~3.5]{Dietzenbacher2021}.
\end{proof}

\section{The additive ray}\label{sec:ray}

\begin{corollary}\label{cor:noncovariance}
Egalitarian stability is not invariant under player-specific additive
transformations, even among exact games with stable cores.
\end{corollary}

\begin{proof}
For the unperturbed game $w$, the first round fixes players $5$ and $6$ at
$1$ through $\{5,6\}$. The second round fixes the other players at $1/2$ and
admits the grand coalition. Thus $w$ is egalitarian stable, with claims
\[
 \left(\tfrac12,\tfrac12,\tfrac12,\tfrac12,1,1\right)
   =\tfrac12(U+V)\in\core(w).
\]
The additively transformed game $v=w+c$ is not egalitarian stable by
Equation~\eqref{eq:claims}.
\end{proof}

The reported vector is one point in a continuous family.  The next result
classifies the simplest ray through it and shows that the failure is not an
isolated tie phenomenon.

\begin{proposition}[an additive phase transition]\label{prop:additive-ray}
For $t\geq0$, put
\[
 c(t)=(t,0,0,0,t,0),\qquad v_t=w+c(t).
\]
Every $v_t$ is monotone, superadditive, exact, and has a stable core.  Its
terminal egalitarian claims are
\[
\widehat\gamma(t)=
\begin{cases}
 (\frac12+\frac34t,\frac12+\frac14t,
  \frac12+\frac34t,\frac12+\frac14t,1+\frac12t,1+\frac12t),
   &0\leq t\leq\frac23,\\[2pt]
 (\frac32t,1-\frac12t,1,1-\frac12t,
  1+\frac12t,1+\frac12t),
   &\frac23\leq t\leq1,\\[2pt]
 (t+\frac12,\frac12,1,\frac12,t+\frac12,\frac32),
   &1\leq t.
\end{cases}
\]
The adjacent formulas agree at the two breakpoints.  Consequently, $v_t$ is
egalitarian stable if and only if $t=0$.
\end{proposition}

\begin{proof}
The preservation claims follow from Lemma~\ref{lem:additive}; monotonicity is
also preserved here because the two nonzero additive coordinates are
nonnegative, while superadditivity is invariant under every additive
transformation.

For the procedural calculation, substitute $v_t(S)=w(S)+c(t)(S)$ into the
equal-share maximizations.  On the three open parameter intervals, the newly
fixed player strata are, respectively,
\[
 56\mid13\mid24,\qquad
 56\mid1\mid234,\qquad
 15\mid36\mid24.
\]
Induction over these strata gives the three displayed affine claim vectors.
At $t=2/3$ and $t=1$ additional maximizing ties occur, but direct substitution
gives the common adjacent formula.  At $t=0$ the grand coalition is admissible.
Finally, since $v_t(N)=4+2t$, the excess of the claim sum over the grand worth is
\[
 \widehat\gamma(t)(N)-v_t(N)=
 \begin{cases}
 t,&0\leq t\leq\frac23,\\
 1-\frac12t,&\frac23\leq t\leq1,\\
 \frac12,&1\leq t.
 \end{cases}
\]
It is zero only at $t=0$.  The package's
\texttt{verification/verify\_additive\_ray.py} checks the affine maximization inequalities
over all three intervals in exact rational arithmetic.
\end{proof}


\section{Discussion}\label{sec:discussion}

The counterexample clarifies the relation among three notions:
\[
\begin{gathered}
 \{\text{large-core games}\}\subsetneq
 \{\text{stable-core games}\},\\
 \{\text{stable-core games}\}\not\subseteq
 \{\text{egalitarian-stable games}\}.
\end{gathered}
\]
The second noninclusion is witnessed within exact, monotone, superadditive
games. The underlying segment geometry is classical; what changes is the
procedural outcome under a nonuniform additive perturbation. This is precisely
where the two stability notions diverge: core stability is invariant under
additive transformations, whereas Corollary~\ref{cor:noncovariance} shows that
egalitarian stability is not.

The mechanism behind stability is local. Any imputation that underpays one of
five pairs is immediately dominated on that pair. Avoiding those pair
dominances forces all but one degree of freedom to agree with the core
parameter. The two triples then point in opposite directions: one dominates
when the residual coordinate is too large, the other when it is too small.
By contrast, after the perturbation the egalitarian procedure combines
different tight coalitions across its two rounds and produces mutually
unaffordable claims. The example therefore shows that external dominance
power of the core does not control the global compatibility of procedural
claims.




\section{Verification and review status}\label{sec:methods}

The six-player counterexample has the direct proof in
Theorem~\ref{thm:main}. The accompanying standard-library Python package is
run with \texttt{python3 verification/run\_checks.py}. It checks the complete
coalition tables, the procedural claims against published positive and
negative controls, the core's affine-line description, and every affine
comparison needed for the additive ray. All finite comparison boundaries are
checked separately. Verification conditions remain active under Python's
optimization flag, and malformed or incomplete coalition data cause failure.

A second exact cross-check enumerates all $20{,}349$ candidate bases of a
$21$-plane arrangement in the efficiency hyperplane, and constructs a core
dominator for each of the $102$ outside-core imputation vertices. It repeats
this check after a specified permutation and additive transformation. The
arrangement argument is recorded in the secondary Appendix~\ref{app:arrangement};
it is not needed for the direct theorem. General polyhedral-core algorithms
are outside this focused manuscript's contribution.

The example was discovered using computational search. Language-model systems
assisted with search, drafting, proof criticism, and verification code. The
present package has received internal mathematical checks and exact
computational reproduction. Independent human specialist review of correctness,
priority, and the correspondence with the original procedural definition is
\emph{pending}; the package makes no claim of completed external review.

\subsection*{Data and code availability}
 the supplied review package contains all
witness data and standard-library verification code; no network access or
proprietary solver is required.

\appendix

\section{A specialized exact cross-check for segment cores}
\label{app:arrangement}

Theorem~\ref{thm:main} has a short direct proof. This appendix records a
general-purpose cross-check used during the audit. It is specialized to a
segment core and complements the general nested-balancedness characterization
of \citet{GrabischSudholter2024}.

Suppose $\core(v)=[A,B]$ and write $d=B-A$ and
$x_t=A+td$ for $0\le t\le1$. Let
\[
 E=\{y:y(N)=v(N)\},\qquad
 I=\{y\in E:y_i\ge v(\{i\})\ \forall i\}.
\]
For a coalition $S$, let
$F_S=\{x\in[A,B]:x(S)=v(S)\}$. Since $x_t(S)-v(S)$ is affine and
nonnegative on $[0,1]$, $F_S$ is empty, $\{A\}$, $\{B\}$, or the entire
segment. A core element that dominates through $S$ must belong to $F_S$,
because core feasibility gives $x(S)\ge v(S)$ while domination requires
$x(S)\le v(S)$.

When $F_S=[A,B]$, put
\[
 P_S=\{i\in S:d_i>0\},\quad
 M_S=\{j\in S:d_j<0\},\quad
 Z_S=\{k\in S:d_k=0\},
\]
and, for $i\in P_S$, $j\in M_S$, define
\[
 h_{ij}(y)=d_j(y_i-A_i)-d_i(y_j-A_j).
\]

\begin{lemma}\label{lem:domformula}
If $F_S=[A,B]$, then a point $y\in I$ is dominated through $S$ by a point of
the segment if and only if
\[
\begin{aligned}
 y_k&<A_k &&(k\in Z_S),\\
 y_i&<B_i &&(i\in P_S),\\
 y_j&<A_j &&(j\in M_S),\\
 h_{ij}(y)&>0 &&(i\in P_S,\ j\in M_S).
\end{aligned}
\]
If $F_S=\{A\}$ or $\{B\}$, the corresponding condition is
$y_i<A_i$ for all $i\in S$ or $y_i<B_i$ for all $i\in S$, respectively.
\end{lemma}

\begin{proof}
For $d_i\ne0$, set $q_i=(y_i-A_i)/d_i$. The inequalities
$x_{t,i}>y_i$ are $t>q_i$ when $d_i>0$ and $t<q_i$ when $d_i<0$; if
$d_i=0$ they are $y_i<A_i$. A parameter $t\in[0,1]$ exists exactly when
\[
 q_i<1\ (i\in P_S),\qquad q_j>0\ (j\in M_S),\qquad
 q_i<q_j\ (i\in P_S,j\in M_S).
\]
These are the displayed inequalities; multiplication by
$d_i>0>d_j$ gives the stated sign for $h_{ij}$. The endpoint cases are
immediate.
\end{proof}

\begin{proposition}[finite segment-core certificate]\label{prop:arrangement}
In the efficiency space $E$, form the hyperplane family
\[
\begin{split}
 {\cal H}={}&
 \{y_i=v(\{i\}),\ y_i=A_i,\ y_i=B_i:i\in N\}\\
 &{}\cup
 \{h_{ij}(y)=0:d_i>0>d_j\}.
\end{split}
\]
The core is stable if and only if every vertex of the ${\cal H}$-arrangement
that lies in $I\setminus[A,B]$ is dominated by a point of $[A,B]$.
\end{proposition}

\begin{proof}
For each $S$, Lemma~\ref{lem:domformula} expresses the set $D_S$ of
imputations dominated through $S$ as an intersection of open halfspaces whose
boundaries belong to ${\cal H}$; the endpoint and empty-face cases have the
same property. Hence $I\setminus D_S$ is a finite union of closed halfspaces
bounded by ${\cal H}$. The undominated set
\[
 U=I\cap\bigcap_{\emptyset\ne S\subsetneq N}(I\setminus D_S)
\]
is therefore a finite union of compact polytopes with facets supported by
${\cal H}$. Also $[A,B]\subseteq U$, because one core point cannot dominate
another.

If $U\setminus[A,B]$ is nonempty, some polytope piece $P$ of this union is not
contained in $[A,B]$. If every vertex of $P$ belonged to the convex segment
$[A,B]$, then $P$, the convex hull of its vertices, would also be contained
in the segment. Thus $P$ has a vertex in $U\setminus[A,B]$. At that vertex,
$n-1$ independent hyperplanes from ${\cal H}$, together with efficiency,
determine the point, so it is an arrangement vertex. The converse is
immediate.
\end{proof}

For the perturbed game, all six entries of $d$ are nonzero, with three
positive and three negative. The individual-rationality floor for each player
duplicates one endpoint-coordinate plane. Thus ${\cal H}$ has
$12+3\cdot3=21$ distinct hyperplanes in the five-dimensional efficiency
space. The exact verifier enumerates all $\binom{21}{5}=20{,}349$ candidate
bases. It obtains $230$ distinct intersections, of which $104$ are
imputations; exactly the two endpoints lie in the core, and each of the
remaining $102$ vertices has a rationally constructed dominator $x_t$.
Proposition~\ref{prop:arrangement} then independently certifies core
stability. The supplement repeats the same check on an equivalent alternate
normalization retained as a regression case.

\begin{thebibliography}{99}

\bibitem[Dietzenbacher and Do\u{g}an(2025)]{DietzenbacherDogan2025}
Dietzenbacher B, Do\u{g}an E (2025) Population monotonicity and egalitarianism.
Int J Game Theory 54:42. \url{https://doi.org/10.1007/s00182-025-00966-y}
Working-paper version: GSBE Research Memoranda 2024/007.
\url{https://doi.org/10.26481/umagsb.2024007}

\bibitem[Gong and Dietzenbacher(2026)]{GongDietzenbacher2026}
Gong D, Dietzenbacher B (2026) Core bound reduced games and consistency.
Int J Game Theory 55:13. \url{https://doi.org/10.1007/s00182-026-00986-2}


\bibitem[Biswas et al.(1999)]{Biswas1999}
Biswas AK, Parthasarathy T, Potters JAM, Voorneveld M (1999) Large cores
and exactness. Games Econ Behav 28:1--12.
\url{https://doi.org/10.1006/game.1998.0686}

\bibitem[Biswas et al.(2001)]{Biswas2001}
Biswas AK, Parthasarathy T, Ravindran G (2001) Stability and largeness of
the core. Games Econ Behav 34:227--237.
\url{https://doi.org/10.1006/game.2000.0804}

\bibitem[Dietzenbacher(2021)]{Dietzenbacher2021}
Dietzenbacher B (2021) Monotonicity and egalitarianism.
Games Econ Behav 127:194--205.
\url{https://doi.org/10.1016/j.geb.2021.03.006}

\bibitem[Dietzenbacher et al.(2017)]{DBH2017}
Dietzenbacher B, Borm P, Hendrickx R (2017) The procedural egalitarian
solution. Games Econ Behav 106:179--187.
\url{https://doi.org/10.1016/j.geb.2017.10.001}

\bibitem[Dutta and Ray(1989)]{DuttaRay89}
Dutta B, Ray D (1989) A concept of egalitarianism under participation
constraints. Econometrica 57:615--635.
\url{https://doi.org/10.2307/1911055}

\bibitem[Est\'evez-Fern\'andez(2012)]{Estevez2012}
Est\'evez-Fern\'andez A (2012) New characterizations for largeness of the
core. Games Econ Behav 76:160--180.
\url{https://doi.org/10.1016/j.geb.2012.04.003}

\bibitem[Grabisch and Sudh\"olter(2024)]{GrabischSudholter2024}
Grabisch M, Sudh\"olter P (2024) Characterization of TU games with stable
cores by nested balancedness. Math Program 203:801--826.
\url{https://doi.org/10.1007/s10107-021-01716-0}

\bibitem[Mauri Masdeu(2017)]{Mauri2017}
Mauri Masdeu L (2017) Essays on egalitarianism-based solution concepts for
cooperative TU-games. Doctoral dissertation, Universitat Rovira i Virgili.
\url{http://hdl.handle.net/10803/460900}

\bibitem[Schmeidler(1972)]{Schmeidler1972}
Schmeidler D (1972) Cores of exact games, I.
J Math Anal Appl 40:214--225.
\url{https://doi.org/10.1016/0022-247X(72)90045-5}

\bibitem[Sharkey(1982)]{Sharkey1982}
Sharkey WW (1982) Cooperative games with large cores.
Int J Game Theory 11:175--182.
\url{https://doi.org/10.1007/BF01755727}

\bibitem[Shellshear and Sudh\"olter(2009)]{Shellshear2009}
Shellshear E, Sudh\"olter P (2009) On core stability, vital coalitions, and
extendability. Games Econ Behav 67:633--644.
\url{https://doi.org/10.1016/j.geb.2009.01.005}

\bibitem[von Neumann and Morgenstern(1944)]{vNM1944}
von Neumann J, Morgenstern O (1944) Theory of Games and Economic Behavior.
Princeton University Press, Princeton

\end{thebibliography}

\end{document}
